% Requires in preamble (consistent with existing macros):
% \newcommand{\src}{\beta}        % volumetric source (birth rate)
% \newcommand{\kr}{a}             % kernel radius
% \newcommand{\cs}{\conc_{s}}     % surface concentration

\subsection{Inclusion of a volumetric source}
Fission-product atoms are generated throughout the kernel volume at a rate set by the
local fission density. Consistent with Assumption~1, this generation is taken to be
uniform over the kernel and constant in time, and enters
\cref{eq:simplified-triso} as a volumetric source $\src$
($\mathrm{mol\ m^{-3}\ s^{-1}}$) on the right-hand side,
%
\begin{equation}
  \pd{\conc}{t}
  = \Diff \left( \pdd{\conc}{r} + \frac{2}{r}\,\pd{\conc}{r} \right) + \src .
  \label{eq:triso-source}
\end{equation}
%
The problem is posed on $0 \le r \le \kr$, where $\kr$ is the kernel radius. In line
with Assumption~2 the outer boundary carries the imposed surface value, and the
solution must remain bounded at the centre despite the $1/r$ curvature term,
%
\begin{equation}
  \conc(\kr,t) = \cs , \qquad
  \left| \conc(0,t) \right| < \infty , \qquad
  \conc(r,0) = \cs .
  \label{eq:bc-ic}
\end{equation}
%
The initial condition states that the kernel starts in equilibrium with its surface,
so all subsequent evolution is driven entirely by the source.

\subsection{Analytical solution}
It is convenient to work with the excess concentration
$v(r,t) = \conc(r,t) - \cs$. Since $\cs$ is constant, every derivative of $v$ equals
the corresponding derivative of $\conc$, so $v$ satisfies \cref{eq:triso-source}
unchanged but with homogeneous data,
%
\begin{equation}
  v(\kr,t) = 0 , \qquad v(r,0) = 0 .
  \label{eq:v-data}
\end{equation}
%
The solution is split into a steady part and a decaying transient,
%
\begin{equation}
  v(r,t) = w(r) + \phi(r,t) .
  \label{eq:split}
\end{equation}

\paragraph{Step 1: steady-state profile.}
The steady part carries the source. Setting $\partial w/\partial t = 0$ in
\cref{eq:triso-source} gives
%
\begin{equation}
  0 = \Diff \left( \frac{\mathrm{d}^{2} w}{\mathrm{d}r^{2}}
    + \frac{2}{r}\,\frac{\mathrm{d}w}{\mathrm{d}r} \right) + \src .
  \label{eq:steady-ode}
\end{equation}
%
Multiplying every term by $r^{2}/\Diff$,
%
\begin{equation}
  r^{2}\,\frac{\mathrm{d}^{2} w}{\mathrm{d}r^{2}}
  + 2r\,\frac{\mathrm{d}w}{\mathrm{d}r}
  = -\frac{\src\, r^{2}}{\Diff} .
  \label{eq:steady-mult}
\end{equation}
%
The left-hand side is an exact derivative. Expanding
$\frac{\mathrm{d}}{\mathrm{d}r}\!\left( r^{2}\,\frac{\mathrm{d}w}{\mathrm{d}r} \right)$
by the product rule,
%
\begin{equation}
  \frac{\mathrm{d}}{\mathrm{d}r}\!\left( r^{2}\,\frac{\mathrm{d}w}{\mathrm{d}r} \right)
  = 2r\,\frac{\mathrm{d}w}{\mathrm{d}r}
  + r^{2}\,\frac{\mathrm{d}^{2} w}{\mathrm{d}r^{2}} ,
\end{equation}
%
which matches \cref{eq:steady-mult} term by term. Hence
%
\begin{equation}
  \frac{\mathrm{d}}{\mathrm{d}r}\!\left( r^{2}\,\frac{\mathrm{d}w}{\mathrm{d}r} \right)
  = -\frac{\src\, r^{2}}{\Diff} .
\end{equation}
%
Integrating once from the origin,
%
\begin{equation}
  r^{2}\,\frac{\mathrm{d}w}{\mathrm{d}r}
  = -\frac{\src}{\Diff} \int r^{2}\,\mathrm{d}r
  = -\frac{\src}{\Diff}\,\frac{r^{3}}{3} + A
  = -\frac{\src\, r^{3}}{3\Diff} + A .
\end{equation}
%
Dividing through by $r^{2}$,
%
\begin{equation}
  \frac{\mathrm{d}w}{\mathrm{d}r}
  = -\frac{\src\, r^{3}}{3\Diff\, r^{2}} + \frac{A}{r^{2}}
  = -\frac{\src\, r}{3\Diff} + \frac{A}{r^{2}} .
\end{equation}
%
As $r \to 0$ the term $A/r^{2}$ diverges unless $A = 0$; boundedness of the flux at
the centre therefore forces $A = 0$. Integrating a second time,
%
\begin{equation}
  w(r) = -\frac{\src}{3\Diff} \int r\,\mathrm{d}r + B
  = -\frac{\src}{3\Diff}\,\frac{r^{2}}{2} + B
  = -\frac{\src\, r^{2}}{6\Diff} + B .
\end{equation}
%
The surface condition $w(\kr) = 0$ fixes the constant,
%
\begin{equation}
  0 = -\frac{\src\, \kr^{2}}{6\Diff} + B
  \quad \Longrightarrow \quad
  B = \frac{\src\, \kr^{2}}{6\Diff} ,
\end{equation}
%
so the steady profile is the parabola
%
\begin{equation}
  w(r) = \frac{\src\, \kr^{2}}{6\Diff} - \frac{\src\, r^{2}}{6\Diff}
  = \frac{\src}{6\Diff}\left( \kr^{2} - r^{2} \right) .
  \label{eq:steady}
\end{equation}
%
The excess concentration is largest at the centre, $w(0) = \src \kr^{2}/(6\Diff)$,
which sets the natural concentration scale of the problem.

\paragraph{Step 2: reduction of the transient to a slab problem.}
The remainder $\phi = v - w$ satisfies the homogeneous (source-free) form of
\cref{eq:triso-source}. Its data follow from \cref{eq:v-data,eq:split}: at the
surface $\phi(\kr,t) = v(\kr,t) - w(\kr) = 0 - 0 = 0$, and initially
$\phi(r,0) = v(r,0) - w(r) = 0 - w(r) = -w(r)$.

The substitution $\psi(r,t) = r\,\phi(r,t)$, equivalently $\phi = \psi/r$, removes
the curvature term. Each derivative of $\phi$ is computed explicitly. The time
derivative, with $r$ held fixed, is
%
\begin{equation}
  \pd{\phi}{t} = \pd{}{t}\!\left( \frac{\psi}{r} \right)
  = \frac{1}{r}\,\pd{\psi}{t} .
\end{equation}
%
The first radial derivative, by the product rule applied to $\psi \cdot r^{-1}$,
%
\begin{equation}
  \pd{\phi}{r}
  = \pd{\psi}{r}\cdot\frac{1}{r} + \psi \cdot \left( -\frac{1}{r^{2}} \right)
  = \frac{1}{r}\,\pd{\psi}{r} - \frac{\psi}{r^{2}} .
\end{equation}
%
Differentiating once more, applying the product rule to each of the two terms,
%
\begin{align}
  \pdd{\phi}{r}
  &= \left[ \frac{1}{r}\,\pdd{\psi}{r}
     + \pd{\psi}{r}\cdot\left( -\frac{1}{r^{2}} \right) \right]
   - \left[ \frac{1}{r^{2}}\,\pd{\psi}{r}
     + \psi \cdot \left( -\frac{2}{r^{3}} \right) \right] \notag \\
  &= \frac{1}{r}\,\pdd{\psi}{r}
   - \frac{1}{r^{2}}\,\pd{\psi}{r}
   - \frac{1}{r^{2}}\,\pd{\psi}{r}
   + \frac{2\psi}{r^{3}}
  = \frac{1}{r}\,\pdd{\psi}{r}
   - \frac{2}{r^{2}}\,\pd{\psi}{r}
   + \frac{2\psi}{r^{3}} .
\end{align}
%
Assembling the diffusion operator and keeping every product before cancelling,
%
\begin{align}
  \pdd{\phi}{r} + \frac{2}{r}\,\pd{\phi}{r}
  &= \frac{1}{r}\,\pdd{\psi}{r}
   - \frac{2}{r^{2}}\,\pd{\psi}{r}
   + \frac{2\psi}{r^{3}}
   + \frac{2}{r}\cdot\frac{1}{r}\,\pd{\psi}{r}
   - \frac{2}{r}\cdot\frac{\psi}{r^{2}} \notag \\
  &= \frac{1}{r}\,\pdd{\psi}{r}
   \underbrace{- \frac{2}{r^{2}}\,\pd{\psi}{r}
   + \frac{2}{r^{2}}\,\pd{\psi}{r}}_{=\,0}
   \underbrace{+ \frac{2\psi}{r^{3}}
   - \frac{2\psi}{r^{3}}}_{=\,0}
  = \frac{1}{r}\,\pdd{\psi}{r} .
\end{align}
%
The homogeneous equation $\phi_t = \Diff(\phi_{rr} + \tfrac{2}{r}\phi_r)$ therefore
becomes $\tfrac{1}{r}\psi_t = \Diff\,\tfrac{1}{r}\psi_{rr}$, and multiplying both
sides by $r$ leaves the one-dimensional heat equation on a slab,
%
\begin{equation}
  \pd{\psi}{t} = \Diff\,\pdd{\psi}{r} , \qquad
  \psi(0,t) = 0 , \qquad \psi(\kr,t) = 0 ,
  \label{eq:slab}
\end{equation}
%
where $\psi(0,t) = 0 \cdot \phi(0,t) = 0$ follows from boundedness of $\phi$ at the
centre, and $\psi(\kr,t) = \kr\,\phi(\kr,t) = 0$ from the surface condition.

\paragraph{Step 3: separation of variables.}
Seek product solutions $\psi(r,t) = \rho(r)\,T(t)$. Substituting into
\cref{eq:slab},
%
\begin{equation}
  \rho(r)\,\frac{\mathrm{d}T}{\mathrm{d}t}
  = \Diff\,T(t)\,\frac{\mathrm{d}^{2}\rho}{\mathrm{d}r^{2}} ,
\end{equation}
%
and dividing both sides by $\Diff\,\rho(r)\,T(t)$,
%
\begin{equation}
  \frac{1}{\Diff\,T}\,\frac{\mathrm{d}T}{\mathrm{d}t}
  = \frac{1}{\rho}\,\frac{\mathrm{d}^{2}\rho}{\mathrm{d}r^{2}}
  = -k^{2} ,
\end{equation}
%
where the separation constant is written $-k^{2}$ because only decaying,
oscillatory-in-space solutions can satisfy two homogeneous boundary conditions. The
spatial problem $\rho'' = -k^{2}\rho$ has the general solution
%
\begin{equation}
  \rho(r) = C_{1}\sin(kr) + C_{2}\cos(kr) .
\end{equation}
%
The condition $\rho(0) = 0$ gives $C_{1}\cdot 0 + C_{2}\cdot 1 = 0$, so $C_{2} = 0$.
The condition $\rho(\kr) = 0$ then requires $\sin(k\kr) = 0$, hence
$k\kr = n\pi$ and
%
\begin{equation}
  k_{n} = \frac{n\pi}{\kr} , \qquad n = 1, 2, 3, \dots
\end{equation}
%
The temporal problem $\mathrm{d}T/\mathrm{d}t = -\Diff k_{n}^{2}\,T$ integrates to
$T_{n}(t) = \exp(-\Diff k_{n}^{2} t)$, with decay rates
%
\begin{equation}
  \lambda_{n} = \Diff k_{n}^{2}
  = \Diff \left( \frac{n\pi}{\kr} \right)^{2}
  = \frac{\Diff\, n^{2}\pi^{2}}{\kr^{2}} .
\end{equation}
%
Superposing all modes,
%
\begin{equation}
  \psi(r,t) = \sum_{n=1}^{\infty} b_{n}
  \sin\!\left( \frac{n\pi r}{\kr} \right)
  \exp\!\left( -\frac{n^{2}\pi^{2}\Diff t}{\kr^{2}} \right) .
  \label{eq:series}
\end{equation}

\paragraph{Step 4: Fourier coefficients from the initial condition.}
At $t = 0$ every exponential equals one, so \cref{eq:series} must reproduce the
initial data $\psi(r,0) = r\,\phi(r,0) = -r\,w(r)$. Writing this out with
\cref{eq:steady},
%
\begin{equation}
  \psi(r,0) = -r \cdot \frac{\src}{6\Diff}\left( \kr^{2} - r^{2} \right)
  = -\frac{\src}{6\Diff}\left( \kr^{2} r - r^{3} \right) .
\end{equation}
%
Orthogonality of the sine modes on $[0,\kr]$ gives the standard coefficient formula
%
\begin{equation}
  b_{n} = \frac{2}{\kr} \int_{0}^{\kr} \psi(r,0)
  \sin\!\left( \frac{n\pi r}{\kr} \right) \mathrm{d}r
  = -\frac{2}{\kr}\cdot\frac{\src}{6\Diff}
  \left[ \kr^{2} I_{1} - I_{3} \right] ,
  \label{eq:bn-def}
\end{equation}
%
where, abbreviating $k = n\pi/\kr$,
%
\begin{equation}
  I_{1} = \int_{0}^{\kr} r \sin(kr)\,\mathrm{d}r , \qquad
  I_{3} = \int_{0}^{\kr} r^{3} \sin(kr)\,\mathrm{d}r .
\end{equation}
%
Both integrals follow from integration by parts, using
$\sin(k\kr) = \sin(n\pi) = 0$ and $\cos(k\kr) = \cos(n\pi) = (-1)^{n}$ at the upper
limit. For $I_{1}$, take $u = r$ and $\mathrm{d}v = \sin(kr)\,\mathrm{d}r$, so that
$v = -\cos(kr)/k$:
%
\begin{align}
  I_{1}
  &= \left[ -\frac{r\cos(kr)}{k} \right]_{0}^{\kr}
   + \frac{1}{k}\int_{0}^{\kr} \cos(kr)\,\mathrm{d}r \notag \\
  &= -\frac{\kr\,(-1)^{n}}{k} + \frac{1}{k}\left[ \frac{\sin(kr)}{k} \right]_{0}^{\kr}
  = -\frac{\kr\,(-1)^{n}}{k} + 0
  = \frac{(-1)^{n+1}\,\kr}{k} .
  \label{eq:I1}
\end{align}
%
For $I_{3}$, take $u = r^{3}$, $v = -\cos(kr)/k$:
%
\begin{equation}
  I_{3}
  = \left[ -\frac{r^{3}\cos(kr)}{k} \right]_{0}^{\kr}
   + \frac{3}{k}\int_{0}^{\kr} r^{2}\cos(kr)\,\mathrm{d}r
  = \frac{(-1)^{n+1}\,\kr^{3}}{k} + \frac{3}{k}\,J_{2} ,
\end{equation}
%
where the remaining integral $J_{2}$ is reduced by one more integration by parts,
with $u = r^{2}$, $\mathrm{d}v = \cos(kr)\,\mathrm{d}r$, $v = \sin(kr)/k$:
%
\begin{equation}
  J_{2} = \int_{0}^{\kr} r^{2}\cos(kr)\,\mathrm{d}r
  = \left[ \frac{r^{2}\sin(kr)}{k} \right]_{0}^{\kr}
   - \frac{2}{k}\int_{0}^{\kr} r\sin(kr)\,\mathrm{d}r
  = 0 - \frac{2}{k}\,I_{1}
  = -\frac{2}{k}\,I_{1} .
\end{equation}
%
Substituting $J_{2}$ back, then $I_{1}$ from \cref{eq:I1},
%
\begin{equation}
  I_{3}
  = \frac{(-1)^{n+1}\,\kr^{3}}{k} + \frac{3}{k}\left( -\frac{2}{k}\,I_{1} \right)
  = \frac{(-1)^{n+1}\,\kr^{3}}{k}
   - \frac{6}{k^{2}}\cdot\frac{(-1)^{n+1}\,\kr}{k}
  = (-1)^{n+1}\left( \frac{\kr^{3}}{k} - \frac{6\kr}{k^{3}} \right) .
\end{equation}
%
Now form the bracket in \cref{eq:bn-def}, keeping both products explicit before
cancelling,
%
\begin{align}
  \kr^{2} I_{1} - I_{3}
  &= \kr^{2}\cdot\frac{(-1)^{n+1}\,\kr}{k}
   - (-1)^{n+1}\left( \frac{\kr^{3}}{k} - \frac{6\kr}{k^{3}} \right) \notag \\
  &= \frac{(-1)^{n+1}\,\kr^{3}}{k}
   - \frac{(-1)^{n+1}\,\kr^{3}}{k}
   + \frac{6\,(-1)^{n+1}\,\kr}{k^{3}}
  = \frac{6\,(-1)^{n+1}\,\kr}{k^{3}} .
\end{align}
%
Restoring $k = n\pi/\kr$, so that $k^{3} = n^{3}\pi^{3}/\kr^{3}$ and
$1/k^{3} = \kr^{3}/(n^{3}\pi^{3})$,
%
\begin{equation}
  \kr^{2} I_{1} - I_{3}
  = 6\,(-1)^{n+1}\,\kr \cdot \frac{\kr^{3}}{n^{3}\pi^{3}}
  = \frac{6\,(-1)^{n+1}\,\kr^{4}}{n^{3}\pi^{3}} .
\end{equation}
%
Substituting into \cref{eq:bn-def} and carrying the prefactors through explicitly,
%
\begin{equation}
  b_{n}
  = -\frac{2}{\kr}\cdot\frac{\src}{6\Diff}\cdot
    \frac{6\,(-1)^{n+1}\,\kr^{4}}{n^{3}\pi^{3}}
  = -\frac{2 \cdot 6 \cdot \src\,(-1)^{n+1}\,\kr^{4}}
        {6\,\Diff\,\kr\, n^{3}\pi^{3}}
  = -\frac{2\src\kr^{3}\,(-1)^{n+1}}{\Diff\, n^{3}\pi^{3}}
  = \frac{2\src\kr^{3}}{\Diff}\,\frac{(-1)^{n}}{n^{3}\pi^{3}} ,
  \label{eq:coeffs}
\end{equation}
%
where the last step used $-(-1)^{n+1} = (-1)^{n}$.

\paragraph{Step 5: full solution.}
The pieces are reassembled through
$\conc = \cs + v = \cs + w + \phi = \cs + w + \psi/r$. Substituting
\cref{eq:steady,eq:series,eq:coeffs},
%
\begin{equation}
  \boxed{\;
  \conc(r,t) = \cs
  + \frac{\src}{6\Diff}\left( \kr^{2} - r^{2} \right)
  + \frac{2 \src \kr^{3}}{\pi^{3} \Diff}
  \sum_{n=1}^{\infty} \frac{(-1)^{n}}{n^{3}}\,
  \frac{1}{r}\sin\!\left( \frac{n\pi r}{\kr} \right)
  \exp\!\left( -\frac{n^{2}\pi^{2}\Diff t}{\kr^{2}} \right) .
  \;}
  \label{eq:full-solution}
\end{equation}
%
Three consistency checks confirm the result. First, the series is regular at the
centre: expanding the sine for small argument,
$\sin(n\pi r/\kr) \approx n\pi r/\kr$, so
$\sin(n\pi r/\kr)/r \to n\pi/\kr$ as $r \to 0$, a finite limit. Second, at $t = 0$
every exponential equals one and the known Fourier identity
$\sum_{n\ge 1} (-1)^{n+1} \sin(nx)/n^{3} = x(\pi^{2}-x^{2})/12$ on $[-\pi,\pi]$,
evaluated at $x = \pi r/\kr$, shows that the series term equals
$-\src(\kr^{2}-r^{2})/(6\Diff)$ exactly, cancelling the parabola and recovering the
uniform initial state $\conc = \cs$. Third, as $t \to \infty$ every mode decays and
the profile relaxes to the steady parabola of \cref{eq:steady}. The slowest mode
($n = 1$) sets the characteristic equilibration time
$\tau = \kr^{2}/(\pi^{2}\Diff)$.

\paragraph{Step 6: release rate.}
The release rate is the outward diffusive flux at the surface multiplied by the
kernel area,
%
\begin{equation}
  R(t) = 4\pi \kr^{2} \times
  \left( -\Diff \left. \pd{\conc}{r} \right|_{r=\kr} \right) .
  \label{eq:release-def}
\end{equation}
%
The radial derivative of \cref{eq:full-solution} is taken term by term. The constant
$\cs$ contributes nothing. The parabola contributes
%
\begin{equation}
  \pd{}{r}\left[ \frac{\src}{6\Diff}\left( \kr^{2} - r^{2} \right) \right]
  = \frac{\src}{6\Diff}\left( 0 - 2r \right)
  = -\frac{2\src\, r}{6\Diff}
  = -\frac{\src\, r}{3\Diff}
  \;\xrightarrow{\;r=\kr\;}\;
  -\frac{\src\, \kr}{3\Diff} .
\end{equation}
%
For the series, the product rule applied to $r^{-1}\sin(n\pi r/\kr)$ gives
%
\begin{equation}
  \pd{}{r}\!\left[ \frac{1}{r}\sin\!\left( \frac{n\pi r}{\kr} \right) \right]
  = \frac{1}{r}\cdot\frac{n\pi}{\kr}\cos\!\left( \frac{n\pi r}{\kr} \right)
  - \frac{1}{r^{2}}\sin\!\left( \frac{n\pi r}{\kr} \right) ,
\end{equation}
%
which at $r = \kr$, using $\cos(n\pi) = (-1)^{n}$ and $\sin(n\pi) = 0$, reduces to
%
\begin{equation}
  \frac{1}{\kr}\cdot\frac{n\pi}{\kr}\,(-1)^{n} - 0
  = \frac{n\pi\,(-1)^{n}}{\kr^{2}} .
\end{equation}
%
Each series term therefore contributes, with all factors written out before
simplifying,
%
\begin{equation}
  \frac{2 \src \kr^{3}}{\pi^{3} \Diff} \cdot \frac{(-1)^{n}}{n^{3}}
  \cdot \frac{n\pi\,(-1)^{n}}{\kr^{2}}\,
  \mathrm{e}^{-\lambda_{n} t}
  = \frac{2 \src \kr^{3}\, n\pi\,(-1)^{2n}}{\pi^{3} \Diff\, n^{3}\, \kr^{2}}\,
  \mathrm{e}^{-\lambda_{n} t}
  = \frac{2 \src \kr}{\pi^{2} \Diff\, n^{2}}\,
  \mathrm{e}^{-\lambda_{n} t} ,
\end{equation}
%
since $(-1)^{2n} = 1$. Collecting the parabolic and series contributions,
%
\begin{equation}
  \left. \pd{\conc}{r} \right|_{r=\kr}
  = -\frac{\src\, \kr}{3\Diff}
  + \frac{2 \src \kr}{\pi^{2} \Diff} \sum_{n=1}^{\infty}
  \frac{1}{n^{2}}\, \mathrm{e}^{-\lambda_{n} t} .
\end{equation}
%
Substituting into \cref{eq:release-def} and multiplying through,
%
\begin{align}
  R(t)
  &= -4\pi \kr^{2} \Diff \left[ -\frac{\src\, \kr}{3\Diff}
   + \frac{2 \src \kr}{\pi^{2} \Diff} \sum_{n=1}^{\infty}
   \frac{\mathrm{e}^{-\lambda_{n} t}}{n^{2}} \right] \notag \\
  &= \frac{4\pi \kr^{2} \Diff\, \src\, \kr}{3\Diff}
   - \frac{4\pi \kr^{2} \Diff \cdot 2 \src \kr}{\pi^{2} \Diff}
   \sum_{n=1}^{\infty} \frac{\mathrm{e}^{-\lambda_{n} t}}{n^{2}}
  = \frac{4}{3}\pi \kr^{3} \src
   - \frac{8 \kr^{3} \src}{\pi}
   \sum_{n=1}^{\infty} \frac{\mathrm{e}^{-\lambda_{n} t}}{n^{2}} .
\end{align}
%
Factoring out the total generation rate $\tfrac{4}{3}\pi \kr^{3} \src$, and noting
that $\tfrac{8}{\pi} \div \tfrac{4\pi}{3} = \tfrac{8 \cdot 3}{\pi \cdot 4\pi}
= \tfrac{6}{\pi^{2}}$,
%
\begin{equation}
  R(t) = \frac{4}{3}\pi \kr^{3} \src
  \left[ 1 - \frac{6}{\pi^{2}} \sum_{n=1}^{\infty} \frac{1}{n^{2}}
  \exp\!\left( -\frac{n^{2}\pi^{2}\Diff t}{\kr^{2}} \right) \right] .
  \label{eq:release}
\end{equation}
%
At $t = 0$ the Basel sum $\sum_{n \ge 1} n^{-2} = \pi^{2}/6$ makes the bracket
$1 - \tfrac{6}{\pi^{2}}\cdot\tfrac{\pi^{2}}{6} = 0$, so no release occurs before a
gradient has developed at the surface. As $t \to \infty$ the sum vanishes and the
release rate rises monotonically toward the total volumetric generation rate, at
which point production and release balance and the kernel inventory saturates. This
is the classical Booth result recovered as a limiting case of the present
formulation.
